% !PRINTED 508
\weberpage{508}
Es seien nämlich $\mu, \nu, \ldots$ die Grade der ganzzahligen Gleichungen, denen (nach 2.) die Functionale $\alpha, \beta, \ldots$ genügen, und $m=\mu\nu\ldots$ Wir verstehen unter $\omega_1, \omega_2, \ldots, \omega_m$ die $m$ Grössen
\[
\alpha^r\beta^s\ldots \qquad r=0,1,\ldots,\mu-1, \quad s=0,1,\ldots,\nu-1, \ldots
\]
Dann können mit Hülfe der Gleichungen $\mu^{\text{ten}}, \nu^{\text{ten}}, \ldots$ Grades, denen die Zahlen $\alpha, \beta, \ldots$ genügen, alle Producte $\alpha^r\beta^s\ldots$, in denen einer der Exponenten $r, s, \ldots$ grösser als $\mu-1, \nu-1, \ldots$ ist, linear in der Form (5) durch $\omega_1, \omega_2, \ldots, \omega_m$ ausgedrückt werden, und dasselbe gilt daher auch von jedem Functional $\omega$ und folglich auch von den Producten $\omega\omega_1, \omega\omega_2, \ldots, \omega\omega_m$. Daraus folgt nach 6., dass $\omega$ ganz ist, wie wir beweisen wollten.

Als speciellen Fall des Satzes 7., aus dem übrigens der allgemeine Satz leicht wieder gefolgert werden kann, heben wir hervor:

8. \gesperrt{Durch Addition, Subtraction und Multiplication ganzer Functionale entstehen immer wieder ganze Functionale.}

Wir haben schon früher (§.\,134) bemerkt, dass man immer einen algebraischen Körper bestimmen kann, der beliebig gegebene algebraische Zahlen enthält. Daraus folgt, dass in dem Satze 7. die $\alpha, \beta, \ldots$ beliebige ganze algebraische Zahlen oder Functionale sein können, und Entsprechendes gilt von dem Satze 8.

Aus der Definition 2. folgt noch nach dem Satze §.\,136, 8.:

9. \gesperrt{Ist $\omega$ ein ganzes Functional des Körpers $\Omega$, so sind auch alle mit $\omega$ conjugirte Functionale ganz.}

Als besondere Folgerung dieses Satzes sei noch erwähnt:

10. \gesperrt{Die absolute Norm eines ganzen Functionales ist eine natürliche ganze Zahl.}

Wir beweisen endlich noch den folgenden Satz:

11. \gesperrt{Wenn ein Functional $\omega$ einer Gleichung von der Form}
\weq{10}{\omega^m + \alpha_1\omega^{m-1} + \cdots + \alpha_m = 0}
\gesperrt{genügt, in der $\alpha_1, \alpha_2, \ldots, \alpha_m$ ganze Functionale in $\Omega$ sind, so ist auch $\omega$ ein ganzes Functional.}
